\subsection{Pairwise Character Geometry and Gram Rigidity}
  \label{subsec:gram-rigidity}

  The character-vanishing condition introduced in the previous section
  imposes strong constraints on the mutual relations between generators.
  These constraints can be expressed in a compact geometric form.

  Let $S=\{s_1,\ldots,s_k\}$ be a generating set of $G$ and let $\rho$
  be an irreducible representation in which each generator of $S$
  satisfies the neutrality condition
  \[
    \chi_\rho(s_i)=0 .
  \]

  The relations between generators can then be encoded in a Gram matrix
  whose entries are determined by representation characters.
  For a pair of generators $s_i,s_j\in S$, we define
  \begin{equation}
    G_{ij}=-\frac{1}{2}\chi_\rho(s_i s_j).
  \end{equation}

  This matrix captures the pairwise geometry of the generating set as
  seen through the representation $\rho$.
  Indeed, when the representation is realised in $SU(2)$,
  each neutral generator corresponds to a traceless element that can be
  written in the form
  \[
    \rho(s_i)=\pm i\,\vec{u}_i\cdot\vec{\sigma},
  \]
  where $\vec{\sigma}$ denotes the vector of Pauli matrices and
  $\vec{u}_i\in S^2$ is a unit vector.

  In this representation one finds
  \[
    \chi_\rho(s_i s_j)=-2\,\vec{u}_i\cdot\vec{u}_j,
  \]
  so that the Gram matrix becomes
  \[
    G_{ij}=\vec{u}_i\cdot\vec{u}_j .
  \]

  The character relations therefore determine the inner products between
  the vectors $\vec{u}_i$ associated with the generators.
  In this sense the geometry of the generating set is fixed by purely
  representation-theoretic data.

  This rigidity has important consequences.
  Once the characters of pairwise products are specified, the relative
  angles between generators cannot be continuously deformed without
  violating the neutrality conditions.
  The admissible configurations of generators therefore form discrete
  families rather than continuous manifolds.

  The Gram formulation provides a convenient bridge between spectral
  admissibility and geometric constraints on generating sets.
  In the next section we exploit this structure by introducing the
  second-moment tensor associated with the vectors $\vec{u}_i$.
  This tensor will allow us to distinguish isotropic configurations
  from anisotropic ones and will lead to a structural exclusion of
  certain candidate groups.
